% ============================================================
%  Chapter 12: Sub-Nucleon Physics and Confinement
%  Source: bounded-phase-space-trajectory.tex, lines 2470--2547
%  Consequences: continuousembed, virtualsubstates, subnucleon
% ============================================================

\epigraph{%
  Quarks are not confined by a force.\\
  They are confined by the structure of the partition.%
}{}

Part~IV of this monograph extends the framework from atomic to nuclear and
sub-nuclear scales.  The two enabling structural results are proved in this
chapter.  The first result---Consequence~\ref{cons:continuousembed}---equips
the discrete partition hierarchy with a unique Riemannian geometry, the
Fisher-information metric, and shows that navigation within it is
logarithmically efficient in the number of leaves.  The second and third
results---Consequences~\ref{cons:virtualsubstates} and
\ref{cons:subnucleon}---show that recursive sub-coordinate decomposition
inside a valid partition state produces intermediate values that are
structurally unobservable, and that sub-nucleon constituents are precisely
such virtual substates.  Quark confinement is therefore not a dynamical
postulate but a geometric consequence of the partition structure.

%------------------------------------------------------------
\section{Continuous Embedding of the Partition Hierarchy}
\label{sec:continuousemb}
%------------------------------------------------------------

The partition hierarchy constructed from Axiom~\ref{ax:bpsl}(iii) is, at
each finite depth, a discrete combinatorial object: a rooted tree whose
leaves are the cells of a given refinement.  The following consequence shows
that this discrete tree has a canonical continuous home.

\begin{consequence}{Continuous Embedding of the Partition Hierarchy}{continuousembed}
The discrete partition hierarchy of Axiom~\ref{ax:bpsl}(iii) with
branching base $b=3$ admits a unique (up to isometry) isometric embedding
into the unit cube $[0,1]^{3}\subset\mathbb{R}^{3}$ equipped with the
Fisher-information Riemannian metric
\begin{equation}
  ds^{2}=\sum_{i=1}^{3}\frac{(dS^{i})^{2}}{S^{i}(1-S^{i})}.
  \label{eq:fishermetric}
\end{equation}
Navigation across the hierarchy (descent from root to leaf) has complexity
$\mathcal{O}(\log_{3}N)$ in this embedding, versus $\Omega(N)$ without the
embedding, where $N$ is the number of partition leaves.
\end{consequence}

\begin{proof}
\emph{Existence.}  Each partition refinement at depth $n$ splits a parent
cell into $b=3$ sub-cells, corresponding to the three spatial dimensions (see
subsection~\ref{subsec:branchingbase} below).  Assign to each depth-$n$ cell
a ternary label $d_{1}d_{2}\cdots d_{n}$ with $d_{k}\in\{0,1,2\}$.  Map
this label to the three ternary fractions
\[
  S^{i}_{n} = \sum_{k=1}^{n} d_{k}^{(i)}\,3^{-k}\in[0,1], \qquad i=1,2,3,
\]
yielding a map $\mathbf{S}\colon \{\text{cells at depth }n\}\to[0,1]^{3}$.
This map is injective: distinct ternary labels produce distinct points.
Its image is dense in $[0,1]^{3}$ as $n\to\infty$; the closure is the entire
cube.

\emph{Uniqueness.}  The Fisher-information metric~\eqref{eq:fishermetric} is
the unique Riemannian metric on $[0,1]^{3}$ invariant under the base-$3$
splitting operations.  This is precisely the statement of Chentsov's
uniqueness theorem applied to the multinomial simplex \citep{chentsov1982}:
among all Riemannian metrics on the space of probability distributions over
three cells, the Fisher metric is the only one invariant under every
sufficient statistic.  Isometry is therefore determined by the metric alone,
and the embedding is unique up to isometry.

\emph{Complexity.}  Navigation from the root to a leaf at depth $n$ follows
a path of length $n=\log_{3}N$ steps, each step choosing one of three
branches.  Without the embedding, finding a specified leaf among $N=3^{n}$
leaves requires in the worst case $\Omega(N)$ comparisons.  The embedding
reduces search to $\mathcal{O}(\log_{3}N)$ steps.
\end{proof}

\subsection{Why the branching base is $b=3$}
\label{subsec:branchingbase}

The choice $b=3$ is fixed by the dimensionality of physical space.  In three
spatial dimensions, the minimum number of sub-cells per parent cell that
preserves the three-axis coordinate structure is three: one independent
refinement per spatial axis.  A branching of two would conflate two axes at
each level and fail to produce independent three-axis resolution.

An equivalent formulation: the ternary splits correspond to the three values
$q\in\{-1,0,+1\}$ of a rank-1 spherical tensor component, so the branching
$b=3$ equals the dimension of the fundamental ($l=1$) representation of
$SO(3)$.  The branching is thus the minimal self-consistent integer for
three-dimensional partitioning under $SO(3)$ symmetry.

\subsection{The Fisher-information metric}
\label{subsec:fishermetric}

The metric~\eqref{eq:fishermetric} on $[0,1]^{3}$ is the pullback of the
standard Fisher metric on the open simplex
$\Delta^{2}=\{(p_{0},p_{1},p_{2}):p_{i}>0,\,\sum p_{i}=1\}$ via the
coordinatewise parametrisation $p_{i}=S^{i}$.  Under this metric, the
geodesic distance between two partition states $\mathbf{S}_{1}$ and
$\mathbf{S}_{2}$ is the Bhattacharyya distance
\[
  d(\mathbf{S}_{1},\mathbf{S}_{2})
  = \arccos\!\left(\sum_{i=1}^{3}\sqrt{S^{i}_{1}S^{i}_{2}}\right),
\]
which measures the statistical distinguishability of the two states.
Navigation along geodesics in the Fisher-metric embedding corresponds to
traversing information-theoretically optimal paths through the hierarchy.

\begin{remark}
The Bhattacharyya distance is related to the quantum fidelity between density
matrices whose diagonals are $(S^{1},S^{2},S^{3})$.  The continuous embedding
therefore connects the discrete combinatorial partition structure to the
standard geometry of quantum state space.
\end{remark}

%------------------------------------------------------------
\section{Virtual Substates and Path Opacity}
\label{sec:virtualsubstates}
%------------------------------------------------------------

Within the unit cube $[0,1]^{3}$, any valid partition state
$\mathbf{S}\in[0,1]^{3}$ can be written as a sum of three sub-coordinates.
The following consequence shows that such decompositions generically produce
intermediate values outside $[0,1]^{3}$ and that these are structurally
invisible.

\begin{consequence}{Virtual Substates and Path Opacity}{virtualsubstates}
Recursive decomposition of a valid global state $\mathbf{S}\in[0,1]^{3}$
into three ternary sub-coordinate triples can produce intermediate
sub-coordinates that lie outside $[0,1]^{3}$.  These intermediate
sub-coordinates are \emph{virtual substates}: they are unobservable by path
opacity, meaning the final global state $\mathbf{S}$ contains no record of
which geodesic through the Fisher-metric embedding was traversed to reach it.
\end{consequence}

\begin{proof}
\emph{Existence of out-of-range sub-coordinates.}  Fix any coordinate
component, say $S^{1}=0.1$.  Decompose as $S^{1}=s_{1}+s_{2}+s_{3}$.  The
system of equations $s_{1}+s_{2}+s_{3}=0.1$ with $s_{i}\in\mathbb{R}$ has
infinitely many solutions, the vast majority of which have at least one
$s_{i}\notin[0,1]$.  For example, $(s_{1},s_{2},s_{3})=(0.5,-0.2,-0.2)$
satisfies the sum but contains negative entries.  Since valid partition cells
require $s_{i}\in[0,1]$, these solutions are not cells of any
$\mathcal{P}_{n}$ and hence are not observable.

\emph{Path opacity.}  Two distinct triples $\{s_{i}\}$ and $\{s'_{i}\}$
satisfying $\sum s_{i}=\sum s'_{i}=S^{1}$ yield the same final state
$\mathbf{S}$.  Any measurement of $\mathbf{S}$ records only $\mathbf{S}$, not
the specific triple that produced it.  The intermediate sub-coordinates are
therefore invisible in principle: no measurement of the final state can
distinguish which path was taken.
\end{proof}

\begin{remark}
Path opacity is not a practical limitation arising from imprecise instruments.
It is structural: the information about which sub-coordinate path was followed
does not exist in the final state.  This is analogous to the fact that the sum
$0.1=0.5+(-0.2)+(-0.2)$ contains no record of the choice of summands; only
the sum $0.1$ is a number.
\end{remark}

\subsection{Examples of virtual substates in recognised physical systems}
\label{subsec:virtualexamples}

Virtual substates arise in several well-known contexts.

\begin{itemize}
\item \textbf{Quarks inside a nucleon.}  Three valence quark ``colours''
  correspond to three sub-coordinate values at the first level of
  sub-decomposition of the nucleon's partition cell.  Their $SU(3)$ colour
  symmetry is the partition symmetry of the ternary split.  By
  Consequence~\ref{cons:virtualsubstates}, individual colours are path-opaque;
  coloured quarks cannot exist in isolation.

\item \textbf{Virtual particles in Feynman diagrams.}  The intermediate
  particles appearing in perturbation-theory diagrams carry four-momenta off
  the mass shell; they are virtual substates in the partition sense.  They
  contribute to amplitudes but are not observable as asymptotic states.

\item \textbf{Pre-decay configurations in radioactive nuclei.}  A metastable
  nucleus between $\alpha$-emission events has a partition sub-structure that
  cannot be directly observed in a single measurement; the specific
  internal configuration before decay is a virtual substate.
\end{itemize}

In each case the virtual substate is both necessary (for the overall
partition dynamics) and unobservable (in isolation).

%------------------------------------------------------------
\section{Sub-Nucleon Structure}
\label{sec:subnucleonstructure}
%------------------------------------------------------------

The partition hierarchy has a coarsest observable scale: the nucleon.  Below
the nucleon radius ($r\lesssim 1$~fm), the following consequence shows that
the hierarchy enters the sub-coordinate regime of
Consequence~\ref{cons:virtualsubstates}.

\begin{consequence}{Sub-Nucleon Constituents as Virtual Substates}{subnucleon}
At spatial resolution below the nucleon radius ($r<1$~fm), the partition
hierarchy enters the sub-coordinate decomposition of
Consequence~\ref{cons:virtualsubstates}.  Putative sub-nucleon constituents
correspond to virtual substates $\{s_{i}\}$; by path opacity they cannot be
isolated as global partition states in principle.
\end{consequence}

\begin{proof}
By Consequence~\ref{cons:continuousembed}, the partition hierarchy embeds in
$[0,1]^{3}$ at a characteristic scale set by the coarsest observable cell.
The coarsest observable global state is the nucleon (Chapter~\ref{ch:nuclear},
Consequence~\ref{cons:nucleonosc}).  At sub-nucleon resolution, the scale
falls below the resolution of the Fisher-metric embedding; the observable
structure at that scale is the sub-coordinate decomposition of the nucleon's
cell.  By Consequence~\ref{cons:virtualsubstates}, the intermediate
sub-coordinates of that decomposition are path-opaque.  Therefore, putative
isolated sub-nucleon constituents are virtual substates and cannot be observed
in isolation.
\end{proof}

\begin{remark}
This deduction simultaneously accounts for two empirical facts about
sub-nucleon structure.  First, deep-inelastic electron-nucleon scattering at
high momentum transfer $Q^{2}$ reveals point-like scattering centres inside
the nucleon \citep{friedman1972}: the virtual substates are real intermediate
objects that deflect the probe.  Second, no free sub-nucleon particle has ever
been observed: the substates are path-opaque and cannot be isolated as global
states.  Both facts are predicted.
\end{remark}

\subsection{Quarks as virtual substates}
\label{subsec:quarks}

The valence quarks of the proton and neutron are the three virtual sub-coordinates
at the first level of ternary decomposition of the nucleon cell.  In the
standard model, quarks carry $SU(3)$ colour charges (red, green, blue); in the
partition framework, these three colours are the three values of the ternary
sub-coordinate $d\in\{0,1,2\}$.  The $SU(3)$ colour group acts as the symmetry
group of the ternary splitting operation, mixing the three sub-coordinates.

\begin{itemize}
\item \textbf{Confinement from path opacity.}  Since the three colour values are
  sub-coordinates of a single nucleon cell, the individual colour is a virtual
  substate.  An isolated coloured quark would require reading off a single
  sub-coordinate independently of the other two, but path opacity prevents
  this: the final cell (the nucleon) contains no record of individual
  sub-coordinate values.  Confinement is thus structural, not dynamical.

\item \textbf{Asymptotic freedom.}  At very high $Q^{2}$ (short distances), the
  probe resolves individual sub-coordinates.  Since these are sub-cells of a
  ternary decomposition, the resolution improves logarithmically with $Q^{2}$
  (the $\mathcal{O}(\log_{3}N)$ complexity of
  Consequence~\ref{cons:continuousembed}).  The weakening of the effective
  sub-coordinate coupling at short distances is the partition-theoretic
  counterpart of asymptotic freedom.

\item \textbf{Fractional charges.}  Each quark carries charge $+\tfrac{2}{3}e$
  or $-\tfrac{1}{3}e$; the proton's net charge $+e$ equals the boundary charge
  of its partition envelope (Consequence~\ref{cons:chargeemergence}).  The
  fractional values are the sub-coordinate fractions of the boundary crossing:
  a sub-coordinate value $d/3$ contributes charge proportional to $d$.  Only
  the boundary of the full nucleon cell has an integer charge assignment.
\end{itemize}

\subsection{Quantitative predictions about sub-nucleon structure}
\label{subsec:deepinelastic}

Deep-inelastic scattering at high $Q^{2}$ measures the momentum distributions
of the virtual substates.  In the partition framework, these distributions are
the probability densities over the ternary sub-coordinate decomposition,
integrated over all allowed geodesics in the Fisher-metric embedding.  The
parton distribution function $q(x)$ measured experimentally gives the
probability that a virtual substate carries momentum fraction $x$ of the
parent nucleon; the ternary sub-coordinate structure predicts, at the
leading-order combinatorial level,
\[
  q(x)\sim x^{-1/2}(1-x)^{2}
\]
for valence quarks at intermediate $x$.  The precise power law receives
corrections from sub-coordinate renormalisation at different scales
(corresponding to QCD evolution equations), which in the partition framework
describe how the sub-coordinate probability distributions transform under
changes of observational resolution.

%------------------------------------------------------------
\bigskip
\begin{tcolorbox}[colback=crownblue!5, colframe=crownblue!60!black,
                  title={\textbf{Chapter 12 Summary}}]
\begin{itemize}[noitemsep]
  \item The discrete partition hierarchy embeds uniquely and isometrically
    into $[0,1]^{3}$ with the Fisher-information metric (Chentsov's theorem),
    reducing navigation complexity from $\Omega(N)$ to $\mathcal{O}(\log_{3}N)$.

  \item Recursive ternary decomposition of a valid partition state produces
    intermediate sub-coordinates outside $[0,1]^{3}$; these are virtual
    substates, invisible by path opacity (Consequence~\ref{cons:virtualsubstates}).

  \item At sub-nucleon resolution the partition hierarchy enters the
    sub-coordinate regime; quarks and other putative sub-nucleon constituents
    are virtual substates and cannot be isolated as global partition states
    (Consequence~\ref{cons:subnucleon}).

  \item Quark confinement is a theorem of partition geometry, not a dynamical
    postulate.  Fractional quark charges are fractional boundary crossings of
    virtual sub-cells.  Asymptotic freedom corresponds to the logarithmic
    navigation efficiency of the Fisher-metric embedding.
\end{itemize}
\end{tcolorbox}
